% =====================================================================
% macros.tex -- notation and editing macros of the manuscript.
% Keep this file small: a macro is worth it only if a symbol appears
% many times or its typesetting may change in one place.
% =====================================================================

% ---- Editing aids (remove or empty \TODO before submission) ---------
\newcommand{\TODO}[1]{\textcolor{red}{[#1]}}
\newcommand{\NOTE}[1]{\textcolor{blue}{[#1]}}

% ---- Generic operators ----------------------------------------------
\newcommand{\dd}{\mathrm{d}}                    % differential
\newcommand{\ee}{\mathrm{e}}                    % Euler's number
\newcommand{\ii}{\mathrm{i}}                    % imaginary unit
\newcommand{\jump}[1]{[\![#1]\!]}               % jump across a point
\newcommand{\T}{\mathsf{T}}                     % transpose
\newcommand{\mat}[1]{\mathbf{#1}}               % matrices and vectors

% ---- Notation decided in step 6.2 (see sections/notation.tex) -------
% Periods and transit times are lower-case t (dimensional) and tau
% (dimensionless, units of L/c_r), like t_a and tau_a; T is kept for
% tensions (T_0, T_d, T_t and the drive tensions T_1, T_2 of DIN 22101).
% Changing one of these macros here changes it everywhere.
\newcommand{\Mtu}{M}                 % take-up mass referred to its pulley
\newcommand{\nstr}{n}                % belt strands that carry the take-up pulley
\newcommand{\Tt}{T_t}                % static belt tension at the take-up
\newcommand{\Wr}{W_r}                % weight of the return strand, g mu_r L
\newcommand{\mw}{m'}                % line density that contributes to weight (DIN 22101 m')
\newcommand{\egrip}{\ee^{\mu\theta}}  % drive factor; E is reserved for Young's modulus
\newcommand{\tper}[1]{t_{#1}}        % natural period of mode #1 (s)
\newcommand{\tauper}[1]{\tau_{#1}}   % the same period in units of L/c_r
\newcommand{\tauB}{\tau_B}           % transit time of strand B, 1 - xi + gamma (units of L/c_r)
\newcommand{\betafive}{\beta_5}      % take-up mass threshold (5 % period change)

% ---- Numbers from the code ------------------------------------------
% paper_numbers.tex defines one macro per number quoted in the text,
% named \Num<Topic><Quantity>, e.g. \NumHarrisonPeriod. Units stay in
% the text: \qty{\NumHarrisonPeriod}{\second}.
